% 第 13 章 ECO Flow
\bichapter{ECO 流程}{ECO Flow}
\label{chap:eco}

工程变更指令（engineering change order，ECO）是对已完成或接近完成的设计所做的增量更改。
可使用 ECO 在无需对整个设计重新综合、布局与布线的情况下，修复功能、时序、噪声与串扰违例；
也可在保持设计性能的同时实现迟到的设计变更。

下列主题描述 Fusion Compiler 工具支持的各种 ECO 流程，以及在这些流程中执行的任务：
\begin{itemize}
  \item 时序或功能变更的通用 ECO 流程（Generic ECO Flow for Timing or Functional Changes）
  \item Freeze Silicon ECO 流程（Freeze Silicon ECO Flow）
  \item Signoff ECO 流程（Signoff ECO Flow）
  \item 增量 Signoff ECO 流程（Incremental Signoff ECO Flow）
  \item ECO Fusion 流程（ECO Fusion Flow）
  \item ECO Fusion 电源完整性流程（ECO Fusion Power Integrity Flow）
  \item 手动实例化备用单元（Manually Instantiating Spare Cells）
  \item 自动添加备用单元（Automatically Adding Spare Cells）
  \item 添加可编程备用单元（Adding Programmable Spare Cells）
  \item 使用 \cmd{eco\_netlist} 命令进行 ECO 更改（Making ECO Changes Using the eco\_netlist Command）
  \item 使用网表编辑命令进行 ECO 更改（Making ECO Changes Using Netlist Editing Commands）
  \item 调整单元尺寸（Resizing Cells）
  \item 在 net 上添加缓冲器（Adding Buffers on Nets）
  \item 在已布线 net 上添加缓冲器（Adding Buffers on Routed Nets）
  \item 优化 net 扇出（Optimizing the Fanout of a Net）
  \item 报告放置 ECO 单元的可用 site（Reporting Available Sites for Placing ECO Cells）
  \item 识别并撤销非最优 ECO 更改（Identifying and Reverting Nonoptimal ECO Changes）
  \item 放置 ECO 单元（Placing ECO Cells）
  \item 将 ECO 单元放置并映射到备用单元（Placing and Mapping ECO Cells to Spare Cells）
  \item 更新 ECO 单元的电源 net（Updating Supply Nets for ECO Cells）
  \item 记录对版图所做的更改（Recording the Changes Made to a Layout）
  \item 执行成组中继器插入与放置的先决条件检查（Performing Prerequisite Check for Group Repeater Insertion and Placement）
  \item 添加一组中继器（Adding a Group of Repeaters）
  \item 查询成组中继器（Querying Group Repeater）
  \item 交换变体单元（Swapping Variant Cell）
  \item 修复多电压违例（Fixing Multivoltage Violations）
\end{itemize}

% ============================================================
\section{时序或功能变更的通用 ECO 流程}
\subsection*{Generic ECO Flow for Timing or Functional Changes}

当可以灵活地添加新单元并移动或删除现有单元时，使用此流程纳入时序或功能 ECO 更改。
若设计尚未 tape out，推荐使用此流程。

无约束（unconstrained）ECO 流程包含下列步骤：
\begin{enumerate}
  \item 用下列方法之一用 ECO 更改更新设计：
  \begin{itemize}
    \item 使用 \cmd{eco\_netlist} 命令，见``使用 \cmd{eco\_netlist} 命令进行 ECO 更改''。
    \item 使用网表编辑 Tcl 命令，见``使用网表编辑命令进行 ECO 更改''。
  \end{itemize}
  \item 使用 \cmd{place\_eco\_cells} 命令更新布局，见``放置 ECO 单元''。
  \item 向 site 阵列的空白处添加填充单元（filler cells），见``插入填充单元（Inserting Filler Cells）''。
  为缩短运行时间，在下列情况下使用 \opt{-post\_eco} 选项：
  \begin{itemize}
    \item 用 \cmd{create\_stdcell\_fillers} 命令插入金属填充单元时，工具将所插入的填充单元标记为 post-ECO 单元。
    \item 用 \cmd{remove\_stdcell\_fillers\_with\_violation} 命令移除存在 DRC 违例的填充单元时，工具仅对 post-ECO 单元执行 DRC 检查。
  \end{itemize}
  \item 使用 \cmd{route\_eco} 命令更新布线，见``执行 ECO 布线（Performing ECO Routing）''，或手动重布受影响的 net。
\end{enumerate}

% ============================================================
\section{Freeze Silicon ECO 流程}
\subsection*{Freeze Silicon ECO Flow}

若单元布局已固定，且只能更改金属与 via 掩模图形，则使用此流程。
若设计已 tape out，并希望避免重新生成整套掩模的费用，推荐使用此流程。

要执行 freeze silicon ECO 流程，块中必须包含备用单元（spare cells）。
可在设计流程的任意时刻用下列方法之一向块添加备用单元：
\begin{itemize}
  \item 手动实例化备用单元，见``手动实例化备用单元''。
  \item 布局后用 \cmd{add\_spare\_cells} 命令自动添加备用单元，见``自动添加备用单元''。
  \item 在芯片收尾（chip finishing）阶段用 \cmd{create\_stdcell\_fillers} 添加可编程备用单元，见``添加可编程备用单元''。
\end{itemize}

Freeze silicon ECO 流程包含下列步骤：
\begin{enumerate}
  \item 将 \appopt{design.eco\_freeze\_silicon\_mode} 应用选项设为 \texttt{true}，以在 freeze silicon 模式下启用 ECO 更改。
  \item 用下列方法之一用 ECO 更改更新设计：
  \begin{itemize}
    \item 使用 \cmd{eco\_netlist} 命令，见``使用 \cmd{eco\_netlist} 命令进行 ECO 更改''。
    \item 使用网表编辑 Tcl 命令，见``使用网表编辑命令进行 ECO 更改''。
  \end{itemize}
  \item 使用 \cmd{check\_freeze\_silicon} 命令分析 ECO 单元到备用单元的映射。
  \item 使用 \cmd{place\_freeze\_silicon} 命令自动将全部 ECO 更改映射到备用单元，或使用 \cmd{map\_freeze\_silicon} 命令将每个 ECO 单元手动映射到特定备用单元，见``将 ECO 单元放置并映射到备用单元''。
  \item 使用 \cmd{route\_eco} 命令更新布线，见``执行 ECO 布线''，或手动重布受影响的 net。
\end{enumerate}

% ============================================================
\section{Signoff ECO 流程}
\subsection*{Signoff ECO Flow}

在 Fusion Compiler 中完成布局布线后，若设计仍有时序或设计规则违例，可在 PrimeTime 中修复这些违例；也可在 PrimeTime 中进行功耗或面积回收（power or area recovery）。

若在 PrimeTime 中对设计做了更改，可生成 ECO 变更列表（change list）文件，并利用 Fusion Compiler 的 ECO 能力将这些更改纳入设计，如图~\ref{fig:fc-signoff-eco} 所示。

\begin{figure}[htbp]
\centering
\fbox{\parbox{0.9\textwidth}{\centering\vspace{1.2cm}
\textit{（原书 Figure 217：Signoff ECO Flow）}\\[0.4em]
\small Fusion Compiler ECO $\leftrightarrow$ ECO change list file $\leftrightarrow$ PrimeTime timing ECO /\\
PrimeTime timing and SI analysis / PrimeTime area and power recovery；\\
StarRC parasitic extraction；Violations? Yes/No
\vspace{1.2cm}}}
\caption{Signoff ECO 流程}
\label{fig:fc-signoff-eco}
\end{figure}

纳入 PrimeTime ECO 更改的步骤如下：
\begin{enumerate}
  \item 通过 source 含有网表编辑命令的 PrimeTime ECO 变更列表 Tcl 文件来更新设计。
  \item 使用 \cmd{place\_eco\_cells} 命令更新布局，示例如下：
\begin{lstlisting}
fc_shell> place_eco_cells -legalize_mode minimum_physical_impact \
   -eco_changed_cells -legalize_only -displacement_threshold 10
\end{lstlisting}
  关于 \cmd{place\_eco\_cells} 的更多信息，见``放置 ECO 单元''。
  \item 向 site 阵列空白处添加填充单元，见``插入填充单元''。
  为缩短运行时间，在下列情况下使用 \opt{-post\_eco} 选项：
  \begin{itemize}
    \item 用 \cmd{create\_stdcell\_fillers} 插入金属填充单元时，工具将所插入填充单元标记为 post-ECO 单元；
    \item 用 \cmd{remove\_stdcell\_fillers\_with\_violation} 移除存在 DRC 违例的填充单元时，工具仅对 post-ECO 单元执行 DRC 检查。
  \end{itemize}
  \item 使用 \cmd{route\_eco} 更新布线，见``执行 ECO 布线''，或手动重布受影响的 net。
\end{enumerate}

% ============================================================
\section{增量 Signoff ECO 流程}
\subsection*{Incremental Signoff ECO Flow}

在 PrimeTime 中做 ECO 更改并经 Fusion Compiler 无约束 ECO 流程纳入后，可能需要在工具间多次迭代才能达到所需 QoR 目标。
为缩短整体周转时间，可用增量 signoff ECO 流程降低每次迭代的运行时间。
该流程由 Fusion Compiler 生成增量设计数据，使你能在 StarRC 中做增量提取，并在 PrimeTime 中做增量时序分析与 ECO，如下图所示。

\begin{figure}[htbp]
\centering
\fbox{\parbox{0.9\textwidth}{\centering\vspace{1.2cm}
\textit{（原书 Figure 218：Incremental Signoff ECO Flow）}\\[0.4em]
\small Fusion Compiler incremental signoff ECO flow $\leftrightarrow$ Incremental design data /\\
ECO change list file $\leftrightarrow$ StarRC incremental parasitic extraction $\leftrightarrow$ Incremental parasitics $\leftrightarrow$\\
PrimeTime incremental timing analysis and ECO $\leftrightarrow$ PrimeTime timing ECO
\vspace{1.2cm}}}
\caption{增量 Signoff ECO 流程}
\label{fig:fc-incr-signoff-eco}
\end{figure}

在 Fusion Compiler 内执行增量 signoff ECO 流程时，使用 \cmd{record\_signoff\_eco\_changes} 命令。
该命令将 PrimeTime ECO 纳入设计库、跟踪对设计所做的全部更改，并生成运行 StarRC 增量提取以及 PrimeTime 增量时序分析与 ECO 所需的增量文件。

Fusion Compiler 增量 signoff ECO 流程包含下列步骤：
\begin{enumerate}
  \item 使用 \cmd{open\_block} 打开设计库。
  \item 使用 \cmd{record\_signoff\_eco\_changes} \opt{-start} \opt{-input} 纳入 PrimeTime ECO 更改并开始跟踪对设计的 ECO 更改，示例如下：
\begin{lstlisting}
fc_shell> record_signoff_eco_changes -start -input pt_eco.tcl
\end{lstlisting}
  \item （可选）使用网表编辑命令对设计做额外时序 ECO 更改。
  \begin{noteBox}
  不要对设计做功能 ECO 更改。若这样做，工具将停止跟踪正在对设计执行的 ECO 更改。
  \end{noteBox}
  \item 使用 \cmd{place\_eco\_cells} 放置 ECO 单元。
  \item 使用 \cmd{route\_eco} 更新布线，见``执行 ECO 布线''，或手动重布受影响的 net。
  \item 使用 \cmd{record\_signoff\_eco\_changes} \opt{-stop} 停止跟踪 ECO 更改并完成增量 signoff ECO 流程，示例如下：
\begin{lstlisting}
fc_shell> record_signoff_eco_changes -stop
\end{lstlisting}
\end{enumerate}

% ============================================================
\section{ECO Fusion 流程}
\subsection*{ECO Fusion Flow}

Fusion Compiler ECO Fusion 流程允许你在 Fusion Compiler 工具内使用 PrimeTime 物理感知 ECO 能力以及 StarRC 提取能力（In-Design signoff extraction）。
可在实现流程的极晚阶段用此流程修复时序与 DRC 违例，并改善面积与功耗 QoR。

Fusion Compiler ECO Fusion 流程包含下列步骤：
\begin{enumerate}
  \item 使用 \cmd{set\_pt\_options} 指定在 Fusion Compiler 内运行 PrimeTime ECO 所需的设置。
  下列示例指定 PrimeTime 可执行文件路径以及分布式处理设置：
\begin{lstlisting}
fc_shell> set_host_options -name pteco_host_option \
   -submit_command "/lsf/bin/bsub -R \"rusage\[mem=\$MEM\]\""
fc_shell> set_pt_options -pt_exec /snps_tools/PT/pt_shell \
   -host_option pteco_host_option
\end{lstlisting}
  若当前工作目录中存在 \texttt{.synopsys\_pt.setup} 文件，Fusion Compiler 会使用其中的信息。
  也可用 \cmd{set\_pt\_options} 的 \opt{-pre\_link\_script} 与 \opt{-post\_link\_script} 指定额外的 PrimeTime 设置。
  更多信息见 \cmd{set\_pt\_options} 的 man page。
  \item 按如下方式设置 StarRC 提取：
  \begin{enumerate}
    \item 确保已启用 StarRC 提取。
    StarRC 提取由 \appopt{extract.starrc\_mode} 应用选项控制。默认值为 \texttt{fusion\_adv}。也可设为 \texttt{in-design} 或 \texttt{none}。
    若要使用原生提取器（native extractor），将该值设为 \texttt{none}。
    \item 使用 \cmd{set\_starrc\_options} \opt{-config} 指定运行 StarRC 提取的配置文件。
  \end{enumerate}
  \item 使用 \cmd{eco\_opt} 命令执行优化。
  默认情况下，\cmd{eco\_opt} 使用 PrimeTime ECO 能力修复时序（setup 与 hold）与 DRC 违例、移除冗余缓冲器，并改善总功耗 QoR；随后用 Fusion Compiler ECO 布局布线能力纳入这些 ECO 更改。
  可用 \opt{-type} 选项控制优化类型。
  下列示例仅用穷尽路径分析（exhaustive path-based analysis）修复时序与 DRC 违例：
\begin{lstlisting}
fc_shell> eco_opt -pba_mode exhaustive -types {timing drc}
\end{lstlisting}
  下列示例改善漏电功耗 QoR，但不纳入 ECO 更改，而是生成包含 ECO 更改的文件：
\begin{lstlisting}
fc_shell> eco_opt -pba_mode exhaustive -types {timing drc} \
   -write_change_file_only
\end{lstlisting}
  更多信息见 \cmd{eco\_opt} 的 man page。
  \item 使用 \cmd{check\_pt\_qor} 分析 PrimeTime QoR。更多信息见该命令的 man page。
\end{enumerate}

\begin{noteBox}
运行 \cmd{eco\_opt} 时，Fusion Compiler 会根据其时序分析设置（如 derating 因子、片上变化设置等）指定 PrimeTime 设置。
因此，为获得最佳收敛，请确保 postroute 阶段的 Fusion Compiler 时序分析设置与 PrimeTime 一致。
可用 \cmd{check\_consistency\_settings} 识别两者设置的差异。
\end{noteBox}

% ============================================================
\section{ECO Fusion 电源完整性流程}
\subsection*{ECO Fusion Power Integrity Flow}

Fusion Compiler ECO Fusion 电源完整性流程允许你使用 RedHawk Fusion 电源轨分析功能识别电源完整性问题，并用 PrimeTime 物理感知 ECO 能力加以修复。

该流程包含下列步骤：
\begin{enumerate}
  \item 设置 \appopt{rail.database} 应用选项，指定保存 RedHawk Fusion 电源轨分析结果的路径。
  \item 使用 \cmd{analyze\_rail} \opt{-voltage\_drop} \texttt{dynamic\_vcd} 或 \cmd{analyze\_rail} \opt{-voltage\_drop} \texttt{dynamic\_vectorless} 执行基于动态向量或不基于向量的压降分析。
  \item 使用 \cmd{set\_pt\_options} 指定在 Fusion Compiler 内运行 PrimeTime ECO 所需的设置。
  \item 使用 \cmd{eco\_opt} \opt{-type} \texttt{power\_integrity} 修复已识别的电源完整性问题。
\end{enumerate}

为修复电源完整性（压降）违例，工具在考虑时序与 DRC QoR 的同时对攻击者单元（aggressor cells）进行缩小尺寸（downsize）。
但该优化要有效，需满足：
\begin{itemize}
  \item 攻击者单元不应处于时序或 DRC 关键；
  \item 单元库中应有更小版本的攻击者单元，且未设置 \texttt{dont\_use} 或 \texttt{dont\_touch} 属性。
\end{itemize}

下列为运行 ECO Fusion 电源完整性流程的示例脚本：
\begin{lstlisting}
set_app_options -name rail.database -value RAIL_DATABASE_BASELINE
analyze_rail -voltage_drop dynamic_vcd -all_nets \
   -switching_activity {VCD ./top.vcd top}
set_pt_options -pt_exec /snps_tools/PT/pt_shell
eco_opt -type power_integrity
\end{lstlisting}

% ============================================================
\section{手动实例化备用单元}
\subsection*{Manually Instantiating Spare Cells}

可用下列方法之一在块中手动实例化备用单元：
\begin{itemize}
  \item 在 Verilog 网表中实例化；
  \item 使用 \cmd{create\_cell}、\cmd{connect\_pins} 等网表编辑命令添加。
\end{itemize}

手动实例化备用单元时，应当：
\begin{itemize}
  \item 在逻辑层次中均匀分布备用单元，以便更好地处理块内任意位置的 ECO 更改；
  \item 将其输入适当地接到电源或地，以防止备用单元产生噪声并消耗功耗。
\end{itemize}

若单元满足下列条件，工具会自动将其识别为备用单元：
\begin{itemize}
  \item 不是仅物理单元（physical-only cell）；
  \item 除时钟 pin 外，所有输入均未连接或接到逻辑常数；备用单元的时钟 pin（若有）可连接到时钟网络；
  \item 所有输出均未连接。
\end{itemize}

若在对块执行物理综合之前实例化备用单元，工具会在后续物理综合步骤中放置并合法化这些备用单元。
但若在已优化、已布局并合法化的块中实例化备用单元，则须用下列步骤放置并合法化：
\begin{enumerate}
  \item 使用 \cmd{spread\_spare\_cells} 分散备用单元。
  默认情况下，该命令在整个 core 区域均匀分布并放置全部备用单元。
  可用下列两种方法之一指定要放置的备用单元：
  \begin{itemize}
    \item 用 \opt{-cells} 放置特定备用单元；
    \item 用 \opt{-voltage\_areas} 放置属于特定电压区域的全部备用单元。
  \end{itemize}
  可如下控制备用单元放置：
  \begin{itemize}
    \item 用 \opt{-boundary} 指定放置备用单元的区域；
    \item 用 \opt{-ignore\_blockage\_types} 忽略特定类型的布局 blockage。默认情况下 \cmd{add\_spare\_cells} 会遵守所有布局 blockage 类型；
    \item 用 \opt{-random\_distribution} 忽略当前单元密度并随机放置备用单元；
    \item 用 \opt{-density\_aware\_ratio} 指定按单元密度分布放置的备用单元百分比，其余随机放置。默认全部按密度分布放置。例如 \opt{-density\_aware\_ratio} \texttt{80} 表示 80\% 按密度、20\% 随机放置。
  \end{itemize}
  \item 使用 \cmd{place\_eco\_cells} \opt{-legalize\_only} \opt{-cells} 合法化备用单元。
\end{enumerate}

% ============================================================
\section{自动添加备用单元}
\subsection*{Automatically Adding Spare Cells}

布局与优化之后，可用下列步骤添加备用单元并合法化：
\begin{enumerate}
  \item 使用 \cmd{add\_spare\_cells} 添加备用单元，并指定下列信息：
  \begin{itemize}
    \item 用 \opt{-cell\_name} 指定备用单元名称前缀；
    \item 用下列两种方法之一指定要插入的备用单元类型与数量：
    \begin{itemize}
      \item 用 \opt{-lib\_cell} 与 \opt{-num\_instances} 指定库单元及每种实例数。例如各插入 250 个 AND2 与 OR2：
\begin{lstlisting}
fc_shell> add_spare_cells -cell_name spare \
   -lib_cell {AND2 OR2} -num_instances 250
\end{lstlisting}
      \item 用 \opt{-num\_cells} 为各库单元指定不同实例数。例如插入 200 个 NAND2 与 150 个 NOR2：
\begin{lstlisting}
fc_shell> add_spare_cells -cell_name spare \
   -num_cells {NAND2 200 NOR 150}
\end{lstlisting}
    \end{itemize}
    \item （可选）用 \opt{-repetitive\_window} 指定在布局区域中重复的放置窗口。
    例如在重复的 $20\times20$ 微米窗口中各添加 15 个 AND2 与 OR2：
\begin{lstlisting}
fc_shell> add_spare_cells -cell_name spare \
  -lib_cell {AND2 OR2} -num_instances 15 \
  -repetitive_window {20 20}
\end{lstlisting}
    在重复的 $25\times20$ 微米窗口中添加 20 个 NAND2 与 15 个 NOR2：
\begin{lstlisting}
fc_shell> add_spare_cells -cell_name spare \
   -num_cells {NAND2 20 NOR 15} -repetitive_window {25 20}
\end{lstlisting}
  \end{itemize}
  默认 \cmd{add\_spare\_cells} 在整个设计中均匀分布备用单元。要在特定范围内分布：
  \begin{itemize}
    \item 层次块：用 \opt{-hier\_cell}。工具在包围该层次块全部单元的矩形区域内分布；
    \item 边界框：用 \opt{-boundary}；
    \item 电压区域：用 \opt{-voltage\_areas}。
  \end{itemize}
  还可进一步控制放置：
  \begin{itemize}
    \item \opt{-ignore\_blockage\_types}：忽略特定类型布局 blockage；
    \item \opt{-random\_distribution}：忽略密度并随机放置；
    \item \opt{-density\_aware\_ratio}：按密度放置的百分比（其余随机）；
    \item \opt{-input\_pin\_connect\_type}：指定备用单元输入 pin 的连接类型，有效值为 \texttt{tie\_low}、\texttt{tie\_high} 或 \texttt{open}。
  \end{itemize}
  \item 使用 \cmd{place\_eco\_cells} \opt{-legalize\_only} \opt{-cells} 合法化备用单元。
\end{enumerate}

% ============================================================
\section{添加可编程备用单元}
\subsection*{Adding Programmable Spare Cells}

若 vendor 提供，Fusion Compiler 支持可编程备用单元（programmable spare cells），亦称门阵列填充单元（gate array filler cells）。
这些单元可通过金属掩模更改进行编程以实现 ECO，从而降低掩模成本并缩短结果交付时间。

插入可编程备用单元的步骤：
\begin{enumerate}
  \item 通过对标准单元与可编程备用单元对应的库单元设置相同的 \texttt{psc\_type\_id} 属性，指定标准单元可映射到的可编程备用单元。
  下列示例指定：
  \begin{itemize}
    \item 名为 BUF1 与 INV1 的标准单元可映射到名为 fill1x 的可编程备用单元（\texttt{psc\_type\_id} 设为 1）；
    \item 名为 NAND2 与 NOR2 的标准单元可映射到名为 fill2x 的可编程备用单元（\texttt{psc\_type\_id} 设为 2）。
  \end{itemize}
\begin{lstlisting}
fc_shell> set_attribute [get_lib_cells fill_lib/fill1x] \
   psc_type_id 1
fc_shell> set_attribute [get_lib_cells stdcell_lib/BUF1] \
   psc_type_id 1
fc_shell> set_attribute [get_lib_cells stdcell_lib/INV1] \
   psc_type_id 1
fc_shell> set_attribute [get_lib_cells fill_lib/fill2x] \
   psc_type_id 2
fc_shell> set_attribute [get_lib_cells stdcell_lib/NAND2] \
   psc_type_id 2
fc_shell> set_attribute [get_lib_cells stdcell_lib/NOR2] \
   psc_type_id 2
\end{lstlisting}
  \item 使用 \cmd{create\_stdcell\_fillers} 插入可编程备用单元。例如插入 fill1x 与 fill2x：
\begin{lstlisting}
fc_shell> create_stdcell_fillers \
        -lib_cells {fill_lib/fill1x fill_lib/fill2x}
\end{lstlisting}
\end{enumerate}

在 ECO 流程中，工具根据 \texttt{psc\_type\_id} 属性设置、单元宽度与电压区域，将 ECO 单元与可编程备用单元交换。
若工具从设计中移除 ECO 单元，可重用对应的可编程备用单元。

% ============================================================
\section{使用 eco\_netlist 命令进行 ECO 更改}
\subsection*{Making ECO Changes Using the eco\_netlist Command}

可使用 \cmd{eco\_netlist} 命令对块进行 ECO 更改。

若使用 freeze silicon ECO 流程，在运行 \cmd{eco\_netlist} 之前，须将 \appopt{design.eco\_freeze\_silicon\_mode} 设为 \texttt{true}。

使用 \cmd{eco\_netlist} 时，用下列两种方法之一进行 ECO 更改：
\begin{itemize}
  \item 用 \opt{-by\_verilog\_file} 指定包含 ECO 更改的 golden Verilog 网表；
  \item 用 \opt{-block} 指定包含 ECO 更改的 golden 块。
  使用 \opt{-block} 时，默认：
  \begin{itemize}
    \item 工具假定 golden 块在当前设计库中；要指定不同设计库，用 \opt{-golden\_lib}；
    \item 工具对当前设计做 ECO 更改；要对不同设计更改，用 \opt{-working\_block}；要用 \opt{-working\_lib} 指定含该设计的设计库。
  \end{itemize}
\end{itemize}

工具将工作设计与输入 Verilog 网表或 golden 设计比较，并生成包含实现功能更改的网表编辑 Tcl 命令的变更文件。
必须用 \opt{-write\_changes} 指定输出文件名。

默认情况下，工具忽略下列差异：
\begin{itemize}
  \item 仅物理单元的差异。要用 \opt{-compare\_physical\_only\_cells} 考虑这些差异；
  \item 时序 ECO 更改，如被调整尺寸的单元或被添加/移除的中继器。要用 \opt{-extract\_timing\_eco\_changes} 在功能 ECO 之外同时考虑时序 ECO；
  \item 电源与地对象的差异。
\end{itemize}

要对工作设计应用 ECO 更改，source \cmd{eco\_netlist} 生成的变更文件，示例如下：
\begin{lstlisting}
fc_shell> eco_netlist -by_verilog_file eco.v \
   -compare_physical_only_cells -write_changes eco_changes.tcl
fc_shell> source eco_changes.tcl
\end{lstlisting}

% ============================================================
\section{使用网表编辑命令进行 ECO 更改}
\subsection*{Making ECO Changes Using Netlist Editing Commands}

若 ECO 更改很少，可用网表编辑命令更新设计。
若使用 freeze silicon ECO 流程，在运行网表编辑命令之前，须将 \appopt{design.eco\_freeze\_silicon\_mode} 设为 \texttt{true}。

网表编辑命令列表见 \textit{Fusion Compiler Data Model User Guide} 中的 Common Design Objects 主题。

\subsection{使用 ECO 脚本进行网表编辑}
\subsubsection*{Using ECO Scripts for Netlist Editing}

Fusion Compiler 支持在设计规划期间使用 ECO 脚本更改网表。可使用：
\begin{itemize}
  \item \cmd{write\_split\_net\_eco}：将物理多扇出 net 的分支上推到顶层；
  \item \cmd{write\_push\_down\_eco}：将标准单元树下推一层；
  \item \cmd{write\_spare\_ports\_eco}：在子块上创建备用端口与 net。
\end{itemize}

更多信息见 \textit{Fusion Compiler Design Planning User Guide} 中的 Generating ECO Scripts for Netlist Editing 主题。

% ============================================================
\section{调整单元尺寸}
\subsection*{Resizing Cells}

可用 \cmd{size\_cell} 调整单元尺寸，示例如下：
\begin{lstlisting}
fc_shell> size_cell U21 -lib_cell AND2X4
\end{lstlisting}

若通过将 \appopt{design.eco\_freeze\_silicon\_mode} 设为 \texttt{true} 启用 freeze silicon 模式，默认情况下 \cmd{size\_cell} 在调整尺寸前会检查在五倍单位 site 高度距离内是否有兼容备用单元。若无兼容备用单元，则不调整该单元。
\begin{itemize}
  \item 用 \opt{-max\_distance\_to\_spare\_cell} 控制该距离；
  \item 用 \opt{-not\_spare\_cell\_aware} 禁用此功能。
\end{itemize}

\subsection{撤销尺寸调整期间所做的更改}
\subsubsection*{Reverting Changes Made During Resizing}

要用 \cmd{revert\_cell\_sizing} 撤销由 \cmd{size\_cell} 所做的 ECO 更改。
用 \opt{-include\_adjacent\_sized\_cells} 可同时包含同一行中与指定单元相邻的已调整尺寸单元。
使用该选项时，默认还撤销指定单元两侧紧邻的已调整尺寸单元；若再使用 \opt{-adjacent\_cell\_distance}，工具会在同一行中递归撤销在指定距离内相邻的已调整尺寸单元。

% ============================================================
\section{在 net 上添加缓冲器}
\subsection*{Adding Buffers on Nets}

要对连接到指定 pin 的 net 添加缓冲器，使用 \cmd{add\_buffer} 命令。
使用时指定下列选项：
\begin{itemize}
  \item \opt{-new\_net\_names}：指定要添加的新 net 名称。添加缓冲器时每个缓冲器须一个 net 名；添加反相器对时每对须两个 net 名。
  \item \opt{-new\_cell\_names}：指定要添加的新单元名称。添加缓冲器时每个缓冲器须一个单元名；添加反相器对时每对须两个单元名。
  \item \opt{-inverter\_pair}：添加反相器对而非缓冲器单元。须提供具有反相输出的库单元。
  \item \opt{-respect\_voltage\_areas}：在电压区域层次内插入缓冲器，使其物理上位于对应布局中。与 \opt{-snap} 互斥。
  \item \opt{-no\_of\_cells}：指定每个 net 要插入的缓冲器单元或反相器对数。
  \item \opt{-object\_list}：指定须缓冲的 net、pin 或端口列表。若其单元已放置，新缓冲器或反相器对会靠近指定 pin 或端口放置。
  \item \opt{-lib\_cell}：指定用作缓冲器的库单元（命名库单元或库单元集合）。此为必需选项，与 \texttt{buffer\_lib\_cell} 互斥。
  \item \opt{-snap}：将新缓冲器放到目标单元旁，对齐到最近 site，并在需要时翻转缓冲器。
\end{itemize}

% ============================================================
\section{在已布线 net 上添加缓冲器}
\subsection*{Adding Buffers on Routed Nets}

要对已完全布线的 net 添加缓冲器或反相器对，使用 \cmd{add\_buffer\_on\_route}，见下列主题：
\begin{itemize}
  \item 指定 net 名称、缓冲器类型及其位置
  \item 控制缓冲器的添加方式
  \item 为多电压设计指定设置
  \item 为 Freeze Silicon ECO 流程指定设置
\end{itemize}

\subsection{指定 net 名称、缓冲器类型及其位置}
\subsubsection*{Specifying the Net Names, Buffers Types, and Their Locations}

使用 \cmd{add\_buffer\_on\_route} 时必须指定：
\begin{itemize}
  \item 要添加缓冲器的 net；
  \item 用下列方法之一指定要添加的缓冲器类型以及在 net 上的位置：
  \begin{itemize}
    \item 按指定配置添加缓冲器
    \item 在指定位置添加缓冲器
    \item 按指定间隔添加缓冲器
    \item 按 net 长度比例的间隔添加缓冲器
    \item 按指定图案在总线（bus）上添加缓冲器
  \end{itemize}
\end{itemize}

\subsubsection{按指定配置添加缓冲器}
\paragraph*{Adding Buffers in a Specified Configuration}

要用 \opt{-user\_specified\_buffers} 按精确配置添加单元及其位置。
该选项一次只能在一条 net 上添加单元；可插入不同类型的缓冲器或反相器对，但不能同时混合两者。

对每个添加的单元，用 \texttt{\{instance\_name library\_cell\_name x y layer\_name...\}} 格式指定：实例名、所用库单元、精确坐标、以及该位置上已有布线图形所在的层。
也可用 \opt{-detect\_layer} 让工具自动检测该位置最近的布线图形，而不指定布线层。

下列示例在 net n22 上添加 ECO1（BUF2，位置 (100,70)，连到 M3）与 ECO2（BUF4，位置 (150,70)，连到 M3）：
\begin{lstlisting}
fc_shell> add_buffer_on_route net22 \
   -user_specified_buffers {ECO1 BUF2 100 70 M3 ECO2 BUF4 150 70 M3}
\end{lstlisting}

\subsubsection{在指定位置添加缓冲器}
\paragraph*{Adding Buffers at Specified Locations}

要在 net 的特定位置添加缓冲器，指定：
\begin{itemize}
  \item 用 \opt{-lib\_cell} 给出可选库单元列表；
  \item 用 \opt{-location} 给出精确位置，格式为 \texttt{\{x1 y1 layer1 ....\}}。也可用 \opt{-detect\_layer} 自动检测最近布线图形。
\end{itemize}
此方法一次只能在一条 net 上添加单元。

下列示例在 net1 的 (100, 200) 处添加 BUF1 并接到 M4：
\begin{lstlisting}
fc_shell> add_buffer_on_route net1 -lib_cell BUF1 \
   -location {100 200 M4}
\end{lstlisting}

默认新单元与 net 的名称前缀为 \texttt{eco\_cell} 与 \texttt{eco\_net}。可用 \opt{-cell\_prefix} 与 \opt{-net\_prefix} 分别指定。

\subsubsection{按指定间隔添加缓冲器}
\paragraph*{Adding Buffers at a Specified Interval}

要对一条或多条 net 按指定间隔添加缓冲器，指定：
\begin{itemize}
  \item \opt{-lib\_cell}：可选库单元列表；
  \item \opt{-first\_distance}：驱动端到第一个缓冲器的距离；
  \item \opt{-last\_distance}：负载 pin 到最后一个缓冲器的距离。
  必须与 \opt{-repeater\_distance} 或 \opt{-repeater\_distance\_length\_ratio} 一起使用；也可与 \opt{-first\_distance} 一起使用。但不能与 \opt{-location}、\opt{-user\_specified\_buffers}、\opt{-user\_specified\_bus\_buffers} 一起使用。
  \begin{noteBox}
  使用 \opt{-last\_distance} 时，\appopt{eco.add\_buffer\_on\_route.min\_distance\_buffer\_to\_load} 应用选项的值被忽略。
  \end{noteBox}
  \item \opt{-repeater\_distance}：缓冲器之间的间隔。
  可选地用 \opt{-scaled\_by\_layer} 为各层指定不同缩放因子，或用 \opt{-scaled\_by\_width} 以层的默认宽度与实际布线宽度之比作为缩放因子。
\end{itemize}

示例：在 n2、n5 上每隔 150 微米添加 BUF1，驱动端到第一个中继器 100 微米：
\begin{lstlisting}
fc_shell> add_buffer_on_route {n2 n5} -lib_cell BUF1 \
  -first_distance 100 -repeater_distance 150
\end{lstlisting}

示例：在 n1 上每隔 100 微米添加 BUF2，首距 50、末距 20：
\begin{lstlisting}
fc_shell> add_buffer_on_route -repeater_distance 100 \
  -first_distance 50 -last_distance 20 net1 -lib_cell BUF2
\end{lstlisting}

\subsubsection{按 net 长度比例的间隔添加缓冲器}
\paragraph*{Adding Buffers at an Interval That is a Ratio of the Net Length}

指定：
\begin{itemize}
  \item \opt{-lib\_cell}；
  \item \opt{-first\_distance\_length\_ratio}：驱动端到第一个缓冲器相对总 net 长度的比例；
  \item \opt{-repeater\_distance\_length\_ratio}：缓冲器间距相对总 net 长度的比例。
\end{itemize}

示例：在 n3、n21、n41 上按总长度 20\% 的间隔添加 BUF1，首距为总长度的 10\%：
\begin{lstlisting}
fc_shell> add_buffer_on_route {n3 n21 n41} -lib_cell mylib/BUF1 \
  -first_distance_length_ratio 0.1 \
  -repeater_distance_length_ratio 0.2
\end{lstlisting}

\subsubsection{按指定图案在总线上添加缓冲器}
\paragraph*{Adding Buffers on a Bus in a Specified Pattern}

步骤如下：
\begin{enumerate}
  \item 使用 \cmd{create\_eco\_bus\_buffer\_pattern} 为总线各 net 定义缓冲器图案，以避免缓冲器成团与重叠。
\end{enumerate}

\begin{table}[htbp]
\centering
\caption{与总线图案相关的命令（原书 Table 72）}
\begin{tabular}{@{}ll@{}}
\toprule
目的 & 命令 \\
\midrule
创建总线图案 & \texttt{create\_eco\_bus\_buffer\_pattern} \\
报告总线图案信息 & \texttt{report\_eco\_bus\_buffer\_patterns} \\
获取缓冲器图案集合 & \texttt{get\_eco\_bus\_buffer\_patterns} \\
移除总线图案 & \texttt{remove\_eco\_bus\_buffer\_patterns} \\
\bottomrule
\end{tabular}
\end{table}

下列示例创建图~\ref{fig:fc-bus-pat-h} 所示水平总线缓冲器图案：第一个缓冲器在最上方 net；从左起以距离 2 交错，每三个缓冲器后重复：
\begin{lstlisting}
fc_shell> create_eco_bus_buffer_pattern -name top_left \
 -first_buffer top -measure_from left -distance 2 -repeat_after 3
\end{lstlisting}

\begin{figure}[htbp]
\centering
\fbox{\parbox{0.85\textwidth}{\centering\vspace{1.0cm}
\textit{（原书 Figure 219：Buffer Patterns for a Horizontal Bus）}\\[0.3em]
\small First buffer；交错距离 2
\vspace{1.0cm}}}
\caption{水平总线的缓冲器图案}
\label{fig:fc-bus-pat-h}
\end{figure}

下列示例创建图~\ref{fig:fc-bus-pat-v} 所示垂直总线图案：第一个缓冲器在最右侧 net；从上起交错距离为 2、3、2，然后重复：
\begin{lstlisting}
fc_shell> create_eco_bus_buffer_pattern -name right_top \
   -first_buffer right -measure_from top \
   -user_specified_distance {2 3 2}
\end{lstlisting}

\begin{figure}[htbp]
\centering
\fbox{\parbox{0.85\textwidth}{\centering\vspace{1.0cm}
\textit{（原书 Figure 220：Buffer Patterns for a Vertical Bus）}\\[0.3em]
\small First buffer；交错距离 2、3、2
\vspace{1.0cm}}}
\caption{垂直总线的缓冲器图案}
\label{fig:fc-bus-pat-v}
\end{figure}

\begin{enumerate}
  \setcounter{enumi}{1}
  \item 使用 \cmd{add\_buffer\_on\_route} \opt{-user\_specified\_bus\_buffers} 在总线 net 上添加缓冲器。
  对每组缓冲器用 \texttt{\{buffer\_pattern\_name library\_cell\_name x y ....\}} 格式指定：图案名、库单元、图案中第一个缓冲器的坐标。
\end{enumerate}

示例：在 data\_A 总线上用图案 BP1 添加一组 BUF2，第一个缓冲器在 (100, 70)：
\begin{lstlisting}
fc_shell> add_buffer_on_route [get_net data_A*] \
   -user_specified_bus_buffers {BP1 BUF2 100 70}
\end{lstlisting}

\subsection{控制缓冲器的添加方式}
\subsubsection*{Controlling How Buffers are Added}

默认情况下，\cmd{add\_buffer\_on\_route}：
\begin{itemize}
  \item 在所有布局 blockage、软宏与硬宏之上添加缓冲器。
  要用 \opt{-dont\_allow\_insertion\_over\_cell} 防止在特定宏上放置；或用 \opt{-respect\_blockages} 防止在所有 blockage 与宏上放置，并可用 \opt{-allow\_insertion\_over\_cell} 列出允许放置的宏。
  \item 在全局布线与详细布线的 net 上都添加缓冲器。
  要用 \opt{-only\_global\_routed\_nets} 仅在尚未分配 track 或详细布线的全局布线 net 上添加。
  \item 在被驱动 pin 的最低公共层次添加缓冲器。
  要用 \opt{-on\_top\_hierarchy} 在尽可能高的层次添加。
  \item 若因 net 的逻辑与物理拓扑差异而需要创建新端口，则不在该布线段上添加缓冲器。
  要用 \opt{-punch\_port} 通过创建新端口在此类布线段上添加。
\end{itemize}

\subsection{为多电压设计指定设置}
\subsubsection*{Specifying Settings for Multivoltage Designs}

对于多电压设计：
\begin{itemize}
  \item 用 \opt{-respect\_voltage\_areas} 确保遵守主/次电压区域，并在物理位于对应电压区域的逻辑层次中添加单元；
  \item 用 \opt{-voltage\_area\_specific\_lib\_cells} 按 \texttt{\{va1 lib\_cell1 va2 lib\_cell2 ...\}} 为不同电压区域指定不同库单元（须同时使用 \opt{-lib\_cell}）；
  \item 用 \opt{-select\_mv\_buffers} 按 \texttt{\{va1 single\_rail\_lib\_cell1 dual\_rail\_lib\_cell1 ...\}} 为不同电压区域指定单轨/双轨库单元（须同时使用 \opt{-lib\_cell}）。未指定的电压区域使用 \opt{-lib\_cell} 中的单元；默认电压区域名为 \texttt{DEFAULT\_VA}；
  \item 用 \opt{-allow\_physical\_feedthrough\_buffer} 允许在指定电压区域的物理 feedthrough net 上放置缓冲器（只能与 \opt{-respect\_voltage\_areas} 或 \opt{-voltage\_area\_specific\_lib\_cells} 一起使用，且须为主电压区域）。之后用 \cmd{set\_eco\_power\_intention} 更新所添加缓冲器的电源 net 与电源域设置；
  \item 用 \opt{-respect\_gas\_station} 在指定电源 net 的 gas station 内添加单元。Gas station 是恒定供电区域，用于对穿越掉电电源域的 net 进行缓冲。用 \opt{-max\_distance\_route\_to\_gas\_station} 指定到 gas station 的最大允许距离。
\end{itemize}

\subsection{为 Freeze Silicon ECO 流程指定设置}
\subsubsection*{Specifying Settings for the Freeze Silicon ECO Flow}

当 \appopt{design.eco\_freeze\_silicon\_mode} 为 \texttt{true} 并在 freeze silicon 模式下运行 \cmd{add\_buffer\_on\_route} 时，工具在添加缓冲器前检查五倍单位 site 高度距离内是否有备用单元；若无则不添加。
\begin{itemize}
  \item 用 \opt{-max\_distance\_to\_spare\_cell} 控制该距离；
  \item 用 \opt{-not\_spare\_cell\_aware} 禁用此功能。
\end{itemize}

% ============================================================
\section{优化 net 扇出}
\subsection*{Optimizing the Fanout of a Net}

在时序 ECO 流程中，可用 \cmd{split\_fanout} 向 net 添加缓冲器并优化其扇出。
必须指定：
\begin{itemize}
  \item 要用 \opt{-net} 按名称或用 \opt{-driver} 按驱动端指定待优化 net；
  \item 用 \opt{-lib\_cell} 指定库单元；
  \item 用下列方法之一指定优化方式：
  \begin{itemize}
    \item \opt{-max\_fanout}：该 net 的最大扇出约束；
    \item \opt{-load}：要缓冲的负载 pin 或端口；可用 \opt{-hierarchy} 指定添加缓冲器的逻辑层次。
  \end{itemize}
\end{itemize}

此外还可以：
\begin{itemize}
  \item 用 \opt{-on\_route} 在现有布线拓扑上添加缓冲器（与 \opt{-load} 合用时不能指定 \opt{-hierarchy}）；
  \item 用 \opt{-max\_distance\_for\_incomplete\_route} 对未完全连接到驱动端或负载端子的已布线 net 拆分扇出（指定端子与不完整布线的最大距离）。若在指定距离内找不到或找到多于一条布线段，工具报错；
  \item 用 \opt{-respect\_blockages} 放置缓冲器时避开布局 blockage 与宏；
  \item 用 \opt{-cell\_prefix} 与 \opt{-net\_prefix} 指定新单元与 net 的名称前缀（默认 \texttt{eco\_cell}/\texttt{eco\_net}）。
\end{itemize}

% ============================================================
\section{报告放置 ECO 单元的可用 site}
\subsection*{Reporting Available Sites for Placing ECO Cells}

可用 \cmd{report\_cell\_feasible\_space} 报告可用于放置 ECO 单元的 site。
默认报告整个块；可用 \opt{-voltage\_areas} 限制到特定电压区域，或用 \opt{-boundary} 限制到特定矩形/多边形区域。

% ============================================================
\section{识别并撤销非最优 ECO 更改}
\subsection*{Identifying and Reverting Nonoptimal ECO Changes}

可用 \cmd{report\_eco\_physical\_changes} 识别由 \cmd{size\_cell}、\cmd{add\_buffer} 与 \cmd{add\_buffer\_on\_route} 所做的 ECO 更改是否会在 ECO 布局与合法化期间导致大位移。

可用 \cmd{revert\_eco\_changes} 撤销这些命令所做的非最优 ECO 更改；用 \opt{-cells} 指定要撤销的单元。

% ============================================================
\section{放置 ECO 单元}
\subsection*{Placing ECO Cells}

对于无约束 ECO 流程，使用 \cmd{place\_eco\_cells} 放置 ECO 单元。
该命令根据连通性与单元相关延迟推导每个 ECO 单元的布局位置，再将每个 ECO 单元合法化到最近的未占用 site。现有单元保持不动，以尽量减小对其布局的影响。

要放置：
\begin{itemize}
  \item 全部未放置单元：用 \opt{-unplaced\_cells}；
  \item 特定单元：用 \opt{-cells}；
  \item 仅因 ECO 操作而更改的单元：用 \opt{-eco\_changed\_cells}。
\end{itemize}

若单元的 \texttt{eco\_change\_status} 属性为下列值之一，则视为已更改：\texttt{create\_cell}、\texttt{change\_link}、\texttt{add\_buffer}、\texttt{size\_cell}。
运行 \cmd{eco\_netlist} 时工具会自动设置该属性；若未使用 \cmd{eco\_netlist} 更新设计，可用 \cmd{set\_attribute} 手动设置。

\subsection{使用 place\_eco\_cells 时控制布局}
\subsubsection*{Controlling Placement When Using the place\_eco\_cells Command}

可如下控制 ECO 单元布局：
\begin{itemize}
  \item \opt{-channel\_aware}：使用宏之间的通道区域进行 ECO 布局（默认避开通道）。
  \item \opt{-fixed\_connection\_net\_weight}：为连接到固定点（如 I/O pad 或宏）的 net 指定权重。默认权重为 1；大于 1 的整数使 ECO 单元更靠近固定点。
  \item 要在 ECO 布局中优先处理特定 net：
  \begin{enumerate}
    \item 用 \cmd{set\_eco\_placement\_net\_weight} 指定 net 权重；
    \item 用 \cmd{place\_eco\_cells} 的 \opt{-honor\_user\_net\_weight} 遵守这些权重。
  \end{enumerate}
  用 \cmd{report\_eco\_placement\_net\_weight} 报告所指定的 net 权重。
  \item 要为 ECO 单元创建虚拟连接并在布局中使用（而非实际连接）：
  \begin{enumerate}
    \item 用 \cmd{create\_virtual\_connection} 创建虚拟连接；
    \item 用 \opt{-use\_virtual\_connection} 在 ECO 布局中使用它们。
  \end{enumerate}
\begin{lstlisting}
fc_shell> create_virtual_connection -name VC_0 \
   -pins {U1/Y ECO17/A}
fc_shell> create_virtual_connection -name VC_1 \
   -pins {U2/Y ECO17/B}
fc_shell> create_virtual_connection -name VC_2 \
   -pins {ECO17/Y D17_OUT} -weight 3
fc_shell> place_eco_cells -cells ECO17 -use_virtual_connection
\end{lstlisting}
  用 \cmd{remove\_virtual\_connections} 移除；用 \cmd{get\_virtual\_connections} 查询。
  \item \opt{-max\_fanout}：忽略连接到 ECO 单元且超过特定阈值的高扇出 net；
  \item \opt{-ignore\_pin\_connection}：忽略 ECO 单元上特定 pin 的连接。
\end{itemize}

\subsection{使用 place\_eco\_cells 时控制合法化}
\subsubsection*{Controlling Legalization When Using the place\_eco\_cells Command}

放置后，\cmd{place\_eco\_cells} 将每个 ECO 单元合法化到最近未占用 site；现有单元保持不动。

可用下列方法控制合法化：
\begin{itemize}
  \item \opt{-no\_legalize}：放置后不合法化；
  \item \opt{-legalize\_only}：仅执行合法化；
  \item \opt{-legalize\_mode}，取值为：
  \begin{itemize}
    \item \texttt{free\_site\_only}（默认）：在空闲 site 上合法化，不移动既有单元。附近无空闲 site 时 ECO 单元可能有大位移；
    \item \texttt{allow\_move\_other\_cells}：通过移动既有单元将 ECO 单元合法化到最近合法位置；
    \item \texttt{minimum\_physical\_impact}：先将可就近合法化到空闲 site 的 ECO 单元合法化；对其余单元再通过移动既有单元合法化。
  \end{itemize}
  \item \opt{-displacement\_threshold}：指定合法化位移阈值；还可与更大的 \opt{-max\_displacement\_threshold} 合用。
\end{itemize}

这些位移阈值与 \opt{-legalize\_mode} 的配合如下：
\begin{itemize}
  \item \texttt{free\_site\_only}：仅合法化在位移阈值内有空闲 site 的 ECO 单元。超过阈值则不合法化，并创建集合 \texttt{epl\_legalizer\_rejected\_cells}。可随后用下式合法化：
\begin{lstlisting}
fc_shell> place_eco_cells \
   -legalize_mode allow_move_other_cells \
   -legalize_only -cells $epl_legalizer_rejected_cells
\end{lstlisting}
  若同时指定 \opt{-max\_displacement\_threshold}，还会创建集合 \texttt{epl\_max\_displacement\_cells}（为前者子集），用于识别位移极大、希望拒绝其 ECO 更改的单元。
  \item \texttt{allow\_move\_other\_cells}：不能指定上述两个位移阈值选项。
  \item \texttt{minimum\_physical\_impact}：先合法化阈值内有空闲 site 的单元；超过阈值的则通过移动既有单元合法化。若还指定 \opt{-max\_displacement\_threshold}，超过该最大阈值的单元不合法化，并放入 \texttt{epl\_max\_displacement\_cells}。
\end{itemize}

\begin{itemize}
  \item \opt{-remove\_filler\_references}：指定合法化期间可移除的填充单元类型。填充单元在非固定布局且与 ECO 单元重叠时被移除；默认合法化时不移除任何填充单元。
  \item 将 \appopt{place.legalize.enable\_advanced\_legalizer} 设为 \texttt{true} 以启用 Advanced Legalizer，并在合法化后移除有违例的填充单元。与 \opt{-remove\_filler\_references} 合用时，仅移除该选项指定的填充单元。此时 \opt{-min\_filler\_distance} 被忽略。
\end{itemize}

\subsection{以最小物理影响（MPI）放置 ECO 单元}
\subsubsection*{Placing ECO Cells With Minimal Physical Impact (MPI)}

为尽量减小对既有布局的扰动，在 \cmd{place\_eco\_cells} 中使用：
\begin{itemize}
  \item \opt{-legalize\_mode} \texttt{minimum\_physical\_impact}
  \item \opt{-max\_displacement\_threshold}
  \item \opt{-remove\_filler\_references}
\end{itemize}

% ============================================================
\section{将 ECO 单元放置并映射到备用单元}
\subsection*{Placing and Mapping ECO Cells to Spare Cells}

对于 freeze silicon ECO 流程，使用 \cmd{place\_freeze\_silicon} 将 ECO 单元放置并映射到备用单元。
默认放置并映射全部 ECO 单元；用 \opt{-cells} 指定特定单元。示例：
\begin{lstlisting}
fc_shell> place_freeze_silicon -cells {ECO1 ECO2 ECO3}
\end{lstlisting}

该命令将每个 ECO 单元放置并映射到最近的匹配备用单元。因 ECO 更改而删除的单元不会从设计中移除，而是转换为备用单元并保留在设计中。

使用 \cmd{place\_freeze\_silicon} 时可以：
\begin{itemize}
  \item 用 \opt{-no\_spare\_cell\_swapping} 将 ECO 单元放到目标备用单元旁但不做映射，以便先分析 QoR；
  \item 用 \opt{-map\_spare\_cells\_only} 仅映射已放到目标备用单元旁的备用单元；
  \item 用 \opt{-write\_map\_file} 生成 ECO 单元到备用单元的映射文件；必要时可编辑后用 \cmd{map\_freeze\_silicon} 映射；
  \item 用 \opt{-min\_filler\_distance} 指定映射到可编程填充单元的 ECO 单元之间的最小距离。
\end{itemize}

\subsection{为可编程备用单元指定映射规则}
\subsubsection*{Specifying Mapping Rules for Programmable Spare Cells}

可用 \cmd{set\_programmable\_spare\_cell\_mapping\_rule}（在运行 \cmd{place\_freeze\_silicon} 之前）定义映射规则，以控制：
\begin{itemize}
  \item 备用单元与 PG net 的重叠；
  \item 多高度备用单元的拆分与单高度备用单元的合并；
  \item 不同类型备用单元的兼容性。
\end{itemize}
用 \cmd{report\_programmable\_spare\_cell\_mapping\_rule} 报告或移除所指定的映射规则。

\subsection{将 ECO 单元映射到特定备用单元}
\subsubsection*{Mapping ECO Cells to Specific Spare Cells}

要用 \cmd{map\_freeze\_silicon} 将 ECO 单元映射到特定备用单元，用 \opt{-eco\_cell} 与 \opt{-spare\_cell} 指定名称。示例：
\begin{lstlisting}
fc_shell> map_freeze_silicon -eco_cell ECO1 -spare_cell spare_1
\end{lstlisting}

也可用映射文件，并通过 \opt{-map\_file} 指定。映射文件格式为：
\begin{lstlisting}
ECO_cell_1   spare_cell_1
ECO_cell_2   spare_cell_2
...      ...
\end{lstlisting}

\subsection{将 ECO 单元映射到逻辑等效备用单元}
\subsubsection*{Mapping ECO Cells to Logically Equivalent Spare Cells}

工具映射 ECO 单元时会查找具有相同库单元名的备用单元；若找不到，该 ECO 单元保持未映射。
要用下列步骤识别并将此类 ECO 单元映射到逻辑等效（LEQ）备用单元：
\begin{enumerate}
  \item 用 \cmd{check\_freeze\_silicon} 识别没有匹配备用单元的 ECO 单元；
  \item 用 \cmd{create\_freeze\_silicon\_leq\_change\_list} 为这些 ECO 单元查找逻辑等效备用单元。指定：
  \begin{itemize}
    \item \opt{-cells}：未映射 ECO 单元名称；
    \item \opt{-output}：输出文件名。该文件为含网表编辑命令的 Tcl 脚本，将每个 ECO 单元替换为逻辑等效备用单元，可以是：
    \begin{itemize}
      \item 具有相同逻辑功能但不同库单元名的单个备用单元；
      \item 最多两个备用单元的组合，以给出相同逻辑功能。
    \end{itemize}
  \end{itemize}
  \item 查看输出 Tcl 文件，确认映射满意，必要时编辑；
  \item 在工具中 source 该 Tcl 文件，用逻辑等效单元替换 ECO 单元。
\end{enumerate}
